% EP24 + EP25. Internal source-relative theorem, 19 October 2026.
% Assumes the full EP20--EP23 trace/quotient infrastructure immediately above.
\section{Automatic core of the genuine terminal phase labels}
\label{sec:ep-automatic-core}

In this section \(\mathcal K_m\) is \emph{terminal}:
no strict same-sort cover of invariant equivalences has
finite relative index. No assertion is made for an arbitrary
terminal unary partition or for safe-looking but
\emph{nonmaximal} phase orientations.

Let \(T=T_m\) denote the finite set of all realised terminal
unary types, and for each \emph{original diversity atom} \(a\)
put \(R_a=\Phi_m(a)\) as in
\eqref{eq:ep-Phi}. These are relations on the finite
label set, and \((t,t)\in R_a\) means that two
\emph{distinct} real vertices of type \(t\) have
old colour \(a\). The old original sort predicates are
included when there are several; our theorem below
starts with precisely one original transitive sort.
No predicate individually naming the \emph{new}
labels is included in the template
\[
   \mathbf T=(T;(R_a)_{a\in L}).
   \tag{C1}\label{eq:ep-template}
\]
An endomorphism of \(\mathbf T\) must preserve all the
\emph{old} atom names. This is the finite core problem
relevant to expansion property, not the exact
equivalence cores used in Part~I.

Call an original atom \(a\) \emph{bipartite} if its
undirected graph on the original \(M\) is bipartite.
Write \(\mathrm{BP}\) for these atoms. Notice that
a bipartite atom need not forbid the monochromatic
triangle for the sole reason that its graph is
bipartite; odd cycles of arbitrary length are
relevant here.

\begin{lemma}[No bipartite atom on a terminal diagonal]
\label{lem:ep-bp-loopless}
For \(a\in\mathrm{BP}\) and any \(t\in T\), one has
\((t,t)\notin R_a\).
\end{lemma}
\begin{proof}
Any restriction of a bipartite graph remains bipartite.
If the current atomic \(a\)-graph within the
homogeneous transitive terminal sort \(t\) were nonempty,
its even-walk equivalence \(Y\) and its connectivity
\(X\) would form invariant, finitely definable
equivalences with \([X:Y]=2\): each connected component
of a nonempty transitive bipartite atomic graph
has exactly two even-walk classes. This is excluded
by terminality. Thus no such edge occurs.
\end{proof}

\begin{lemma}[Original functional involutions]
\label{lem:ep-original-functional}
Let \(\mathscr H\) be the group of all
\(\operatorname{Aut}(M)\)-equivariant permutations
of the underlying original transitive set \(M\).
It is a finite elementary abelian \(2\)-group.
For nonidentity \(h\in\mathscr H\), its graph is
one original bipartite atomic relation. Each
\(h\) acts as a permutation on \emph{every}
prefix-type set \(T_j\); a nonidentity \(h\)
acts freely on \(T\).
\end{lemma}
\begin{proof}
By transitivity, the graph of \(h\) is one
old orbital (or logical identity). Old same-sort
symmetry gives \(h^2=1\). There are finitely
many orbitals, so \(\mathscr H\) is a finite group
of exponent two and is elementary abelian.
A nonidentity graph is a perfect matching
and is bipartite. By
\cref{lem:ep-phi-support}, the support
\(R=\Phi_j(\operatorname{graph}(h))\)
is a relation on the finite prefix-type
set \(T_j\), is symmetric because the
original orbital is symmetric, and satisfies
\[
                    R;R=\Id_{T_j}.
\]
It is also total on \(T_j\): every
original point has an \(h\)-image and
each prefix type is realized.
For clarity, these properties FORCE
\(R\) to be the graph of a unique
involution. Indeed, if a type \(p\)
had two neighbours \(q,r\) in \(R\),
symmetry would put \((q,r)\) in
\(R;R=\Id_{T_j}\), so \(q=r\).
Totality gives exactly one neighbour
of every type; symmetry makes the
resulting permutation involutive.
Thus EVERY original equivariant
permutation \(h\) acts on EVERY
prefix-type set, although \(h\) need
not itself preserve the individual
expanded unary predicates.
At terminality its action has no fixed
type by \cref{lem:ep-bp-loopless}.
\end{proof}

\begin{lemma}[Lifting bipartite quotient separators]
\label{lem:ep-lift-separator}
Let \(E>\Id\) be original full exact, and let
\(s,t\in T\) belong to \emph{different}
\(\Phi_m(E)\)-classes \(C,D\). Suppose the
quotient terminal tower on \(M/E\) has an
original quotient bipartite atom \(b\)
occurring between \(C,D\). Then an
\emph{original} bipartite atom \(a\) of \(M\)
occurs between \(s,t\).
\end{lemma}
\begin{proof}
Choose an actual \(b\)-pair of \(E\)-classes \(X,Y\)
of occupancy types \(C,D\). By
\cref{lem:ep-phi-support}, \(X\) contains an
\(s\)-point and \(Y\) a \(t\)-point. Choose them.
Their old atom \(a\) belongs to the
\(E\)-saturated pullback of \(b\).
A bipartition of the quotient \(b\)-graph
pulls back to a bipartition of the old
\(a\)-graph, because every old \(a\)-edge is
contained in that same saturated quotient
orbital. Thus \(a\in\mathrm{BP}\) and
\((s,t)\in R_a\). No distinguished raw
representatives or common ordering have been
used.
\end{proof}

\begin{theorem}[Automatic complete separation and core]
\label{thm:ep-automatic-core}
For every genuinely generated terminal tower
\eqref{eq:ep-tower} over the original one-sort
symmetric AP+homomorphic-CCO source, one has
\[
 \boxed{\quad s\ne t\ \Longrightarrow\
    \exists a\in\mathrm{BP}\ 
       (s,t)\in R_a .\quad}
 \tag{C2}\label{eq:ep-bp-separation}
\]
Consequently the graph of terminal types joined
by a common globally loopless old atom is
\emph{complete} and every endomorphism of
\(\mathbf T\) is an automorphism. In particular,
\(\mathbf T\) is a \emph{core}.
\end{theorem}
\begin{proof}
We prove \eqref{eq:ep-bp-separation} by strong
induction on the number \(k\) of original
diversity orbitals. The induction category includes
all one-sort homogeneous symmetric AP+homomorphic-CCO
sources, including finite and non-SAP structures:
an exact quotient of an SAP source need not have SAP.

If \(k=0\), the original one-sort complete structure
has at most one point and there is at most one realised
terminal label. A zero-step tower likewise has
one label. Assume \(k>0\) and a nonempty tower.
Let \(G_1,\ldots,G_m\) be its \emph{original}
canonical gates, one per connected safe step, from
\cref{lem:ep-single-old-gate}.

\emph{Case 1: \(G_i>\Id\) for every step.}
By \eqref{eq:ep-common-kernel}, for distinct
\(s,t\in T\) some \(i\) satisfies
\[
       [s]_{\Phi_m(G_i)}
            \ne [t]_{\Phi_m(G_i)} .
\]
The projected tower on \(M/G_i\) is genuine
and terminal by
\cref{thm:ep-quotient-tower}. Since \(G_i>\Id\)
its original diversity rank is strictly below \(k\)
(\cref{rem:ep-trace-audit}). The induction
hypothesis supplies a bipartite quotient
atom joining these two quotient labels.
Apply \cref{lem:ep-lift-separator} to obtain
an original bipartite separator between \(s,t\).

\emph{Case 2: some \(G_i=\Id\).}
The original gate \(\Id\) is maximal for one
original parity pivot \(c\). Its even relation
\(Y_c\) is contained in \(\Id\); thus
\(Y_c=\Id\). The symmetric nonempty
old orbital \(c\) is serial on the
original transitive sort, so
\(\Id\subseteq c;c\). But every
two-edge \(c\)-walk is even, hence
\(c;c\subseteq Y_c=\Id\). Therefore
\(c;c=\Id\). A symmetric serial
relation whose square is identity
is the graph of a unique nontrivial
fixed-point-free involution
\(u\in\mathscr H\). The selected
character takes this involution to one.
Original maximality says that the graph of \(u\)
is contained in \emph{every} original full
exact equivalence strictly above \(\Id\).

The finite meet
\[
 F=\bigcap\{E\in\operatorname{ExEq}(M):
                                      E>\Id\}
 \tag{C3}\label{eq:ep-least-F}
\]
is therefore an original full exact equivalence
\emph{strictly} above \(\Id\), and the least
such one. Put
\[
 \mathscr H_F=
    \{h\in\mathscr H:\operatorname{graph}(h)\subseteq F\}.
 \tag{C4}\label{eq:ep-HF}
\]
It is a nontrivial subgroup.
Every original gate is now either \(\Id\) or
contains \(F\).

We claim that, \emph{for every prefix \(j\)},
\[
 \mathscr H_F\text{ acts transitively on each }
     \Phi_j(F)\text{-class of realised prefix types}.
 \tag{C5}\label{eq:ep-HF-transitive}
\]
Induct on the prefix length. Initially there is
one original type, so the assertion is immediate.

Suppose it holds before a step with original
gate \(G>\Id\). Then \(F\le G\).
By \eqref{eq:ep-graph-offsets}, each active
old \(\Phi(F)\)-class acquires two new
occupancy classes, each projecting \emph{bijectively}
to the old class, with objectwise character offsets.
The group \(\mathscr H_F\) preserves both
new classes, since its graphs are contained in
\(F\). Projection is equivariant, so the
transitivity of the old class transports to
each new class. An inactive class is unchanged.

Now suppose the step has original gate \(G=\Id\).
By \cref{lem:ep-single-old-gate}, its active
parents form one orbit of the full original
functional group \(\mathscr H\).
Since \(\mathscr H\) is abelian,
\(\mathscr H_F\) is normal, and each
\(\mathscr H_F\)-orbit is either wholly active
or wholly inactive. By induction each
\(\Phi(F)\)-class is one such orbit.
For an active parent \(p\), the selected
odd functional involution is an original
element of \(\mathscr H_F\). To check
this independently of a chosen label
for a groupoid object, fix the source
parent \(p_*\) and let \(u\) be the
original functional atom selected there.
Its graph belongs to \(F\), by maximality
of the old identity gate. It fixes the
source parent type: the chosen
same-parent odd pivot has a nonempty
trace \(p_*^2\) and
\(\Phi_j(u)\) is a permutation, so
\(u(p_*)=p_*\).

For any other active parent
\(p=h(p_*)\), where \(h\in\mathscr H\)
because the FULL old identity-gate
groupoid is one original functional
\(\mathscr H\)-orbit, commutativity
of the elementary abelian group gives
\[
                 u(p)=u h(p_*)=
                      h u(p_*)=p.
\]
The transported groupoid odd element
at \(p\) is \(huh^{-1}=u\);
its character still equals one, so
the SAME original involution \(u\in
\mathscr H_F\) exchanges both realized
new children over EVERY active parent.
This requires only the action of original
equivariant relations on current type
supports, not an assertion that \(h\)
is itself an automorphism of the chosen
marked expansion.
The old \(\mathscr H_F\)-transitivity moves
any active parent \(p\) to any other one
\(q\) in the same occupancy class.
Then the SAME \(u\in\mathscr H_F\)
exchanges the two children over \(q\).
Thus \(\mathscr H_F\) acts transitively on
the \emph{single} new \(\Phi(F)\)-class
consisting of all the children over that
old class. Inactive classes remain unchanged.
This establishes \eqref{eq:ep-HF-transitive}
for every \(j\).

Return to arbitrary distinct terminal \(s,t\).
If their \(\Phi_m(F)\)-classes differ,
quotient by \(F>\Id\), apply the diversity-rank
induction hypothesis to the genuine terminal
tower on \(M/F\), and lift the bipartite
separator using \cref{lem:ep-lift-separator}.
If their classes agree, by
\eqref{eq:ep-HF-transitive} there is
\(h\in\mathscr H_F\) with \(h(s)=t\).
Since \(s\ne t\), \(h\) is nonidentity.
Its graph is an original bipartite atom
by \cref{lem:ep-original-functional},
and its type image witnesses
\((s,t)\in R_h\). This also separates \(s,t\).

This finishes the induction. Finally, if
an endomorphism \(f\) identified \(s\ne t\),
an atom \(a\in\mathrm{BP}\) separating them
would give the forbidden diagonal
\((f(s),f(s))\in R_a\), contradicting
\cref{lem:ep-bp-loopless}.
Thus all endomorphisms are injective and,
on a finite relational structure, automorphisms.
\end{proof}

\begin{remark}[Proof boundaries and negative control]
\label{rem:ep-core-scope}
This theorem is stronger than core-ness alone;
it gives an original-bipartite-atom separator
for \emph{every pair} of terminal labels.
No intermediate step invokes Ramsey,
expansion property or an existing core theorem,
and the induction is on original
\emph{diversity rank}, not on tower length.

It is false if maximality of the chosen exact
complements is omitted, even while reachable
parity, full character transport, homogeneity,
SAP and CCO are retained. The explicit
four-atom source \(I\times\mathbb F_2\times\mathbb N\)
from the EP18 proof record has a safe first
parity orientation using the \emph{nonmaximal}
clone gate, followed by a maximal split;
its four-label table has the proper retraction
\((\lambda,\mu)\mapsto(\mu,\mu)\).
Choosing the actual first maximal gate
instead gives two terminal labels and a core.

Neither this theorem nor the shorter terminal
diagonal-palette lemma identifies a unique
preferred safe tower, proves a least vocabulary,
or extends without work to several \emph{original}
unary sorts. The source-level exact/quotient/phase
inputs require independent review before
formal verification or publication.
\end{remark}
